\section{Feature Engineering}
Raw telemetry attributes alone are insufficient to represent compound thermal and capacity stress on substation transformers. We engineered 7 domain-specific indicators, expanding the feature space from 27 raw variables to 28 model inputs:

\begin{enumerate}
    \item \textbf{Load Utilization Ratio ($U_{\text{load}}$):} Quantifies transformer loading relative to rated capacity:
    \begin{equation}
    U_{\text{load}} = \frac{\text{Transformer Load}}{\text{Transformer Capacity} + 10^{-5}}
    \end{equation}
    \item \textbf{Demand Utilization Ratio ($U_{\text{demand}}$):} Expresses aggregate electricity demand relative to transformer rating:
    \begin{equation}
    U_{\text{demand}} = \frac{\text{Electricity Demand}}{\text{Transformer Capacity} + 10^{-5}}
    \end{equation}
    \item \textbf{Renewable Generation Ratio ($R_{\text{ratio}}$):} Measures renewable generation contribution relative to local demand:
    \begin{equation}
    R_{\text{ratio}} = \frac{\text{Renewable Generation}}{\text{Electricity Demand} + 10^{-5}}
    \end{equation}
    \item \textbf{Temperature-Humidity Index ($\text{THI}$):} Evaluates environmental heat stress accelerating thermal insulation degradation:
    \begin{equation}
    \text{THI} = T + 0.55 \times (1 - H/100) \times (T - 14.5)
    \end{equation}
    where $T$ denotes ambient temperature ($^{\circ}\text{C}$) and $H$ denotes relative humidity (\%).
    \item \textbf{Thermal-Wind Interaction:} Models convective cooling effects on exposed substation equipment via the product $T \times \text{Wind Speed}$.
    \item \textbf{Temporal Flags:} Binary indicators \texttt{is\_peak\_hour} (18:00--22:00) and \texttt{is\_weekend} capture cyclical human consumption behaviors.
\end{enumerate}
